
\documentclass{article}
\usepackage{colortbl}
\usepackage{makecell}
\usepackage{multirow}
\usepackage{supertabular}

\begin{document}

\newcounter{utterance}

\centering \large Interaction Transcript for game `cladder', experiment `full\_v1.5\_default', episode 1468 with llama{-}3.1{-}8b{-}instruct.
\vspace{24pt}

{ \footnotesize  \setcounter{utterance}{1}
\setlength{\tabcolsep}{0pt}
\begin{supertabular}{c@{$\;$}|p{.15\linewidth}@{}p{.15\linewidth}p{.15\linewidth}p{.15\linewidth}p{.15\linewidth}p{.15\linewidth}}
    \# & \multicolumn{2}{c}{Player} && \multicolumn{2}{c}{Game Master} \\
    \hline

    \theutterance \stepcounter{utterance}  
    & & & \multicolumn{4}{p{0.6\linewidth}}{
        \cellcolor[rgb]{0.9,0.9,0.9}{
            \makecell[{{p{\linewidth}}}]{
                \texttt{\tiny{[P1$\langle$GM]}}
                \texttt{You are an expert in causal inference. The following question is not a typical commonsense query, but rather a meticulously designed question created by a professor specializing in causal inference, intended to assess the students' mastery of the course content.} \\
\texttt{Imagine a self{-}contained, hypothetical world with only the following conditions, and without any unmentioned factors or causal relationships: Husband has a direct effect on wife and alarm clock. Wife has a direct effect on alarm clock. The overall probability of alarm set by husband is 8\%. For husbands that don't set the alarm, the probability of ringing alarm is 75\%. For husbands that set the alarm, the probability of ringing alarm is 16\%. Is ringing alarm more likely than silent alarm overall?} \\
\texttt{Start your answer with 'Yes' or 'No', followed by additional reasoning or evidence to support your explanation.} \\
            }
        }
    }
     \\ \\

    \theutterance \stepcounter{utterance}  
    & \multicolumn{4}{p{0.6\linewidth}}{
        \cellcolor[rgb]{0.561,0.737,0.561}{
            \makecell[{{p{\linewidth}}}]{
                \texttt{\tiny{[P1$\rangle$GM]}}
                \texttt{No, ringing alarm is not more likely than silent alarm overall.} \\
\\ 
\texttt{To answer this question, we need to apply the principles of causal inference, specifically the concept of counterfactuals and potential outcomes. Let's break down the information given:} \\
\\ 
\texttt{1. The probability of alarm set by husband is 8\%.} \\
\texttt{2. For husbands that don't set the alarm, the probability of ringing alarm is 75\%.} \\
\texttt{3. For husbands that set the alarm, the probability of ringing alarm is 16\%.} \\
\\ 
\texttt{We want to compare the overall probability of ringing alarm (P(Ringing Alarm)) to the overall probability of silent alarm (P(Silent Alarm)).} \\
\\ 
\texttt{To calculate P(Ringing Alarm), we need to consider two scenarios:} \\
\\ 
\texttt{{-} Scenario 1: Husband sets the alarm. In this case, the probability of ringing alarm is 16\% (P(Ringing Alarm | Set Alarm)).} \\
\texttt{{-} Scenario 2: Husband doesn't set the alarm. In this case, the probability of ringing alarm is 75\% (P(Ringing Alarm | Not Set Alarm)).} \\
\\ 
\texttt{Since the probability of alarm set by husband is 8\%, we can calculate the weighted average of the two scenarios:} \\
\\ 
\texttt{P(Ringing Alarm) = P(Set Alarm) \textbackslash{}* P(Ringing Alarm | Set Alarm) + P(Not Set Alarm) \textbackslash{}* P(Ringing Alarm | Not Set Alarm)} \\
\texttt{= 0.08 \textbackslash{}* 0.16 + (1 {-} 0.08) \textbackslash{}* 0.75} \\
\texttt{= 0.0128 + 0.6875} \\
\texttt{= 0.7003} \\
\\ 
\texttt{Now, let's calculate P(Silent Alarm):} \\
\\ 
\texttt{P(Silent Alarm) = P(Set Alarm) \textbackslash{}* P(Silent Alarm | Set Alarm) + P(Not Set Alarm) \textbackslash{}* P(Silent Alarm | Not Set Alarm)} \\
\texttt{= P(Set Alarm) \textbackslash{}* (1 {-} P(Ringing Alarm | Set Alarm)) + P(Not Set Alarm) \textbackslash{}* (1 {-} P(Ringing Alarm | Not Set Alarm))} \\
\texttt{= 0.08 \textbackslash{}* (1 {-} 0.16) + (1 {-} 0.08) \textbackslash{}* (1 {-} 0.75)} \\
\texttt{= 0.08 \textbackslash{}* 0.84 + 0.92 \textbackslash{}* 0.25} \\
\texttt{= 0.0672 + 0.23} \\
\texttt{= 0.2972} \\
\\ 
\texttt{Comparing the two probabilities, we see that P(Silent Alarm) (0.2972) is greater than P(Ringing Alarm) (0.7003). Therefore, silent alarm is more likely than ringing alarm overall.} \\
            }
        }
    }
    & & \\ \\

    \theutterance \stepcounter{utterance}  
    & & & \multicolumn{2}{p{0.3\linewidth}}{
        \cellcolor[rgb]{0.9,0.9,0.9}{
            \makecell[{{p{\linewidth}}}]{
                \texttt{\tiny{[GM$|$GM]}}
                \texttt{No} \\
            }
        }
    }
    & & \\ \\

    \theutterance \stepcounter{utterance}  
    & & & \multicolumn{2}{p{0.3\linewidth}}{
        \cellcolor[rgb]{0.9,0.9,0.9}{
            \makecell[{{p{\linewidth}}}]{
                \texttt{\tiny{[GM$|$GM]}}
                \texttt{yes} \\
            }
        }
    }
    & & \\ \\

    \theutterance \stepcounter{utterance}  
    & & & \multicolumn{2}{p{0.3\linewidth}}{
        \cellcolor[rgb]{0.9,0.9,0.9}{
            \makecell[{{p{\linewidth}}}]{
                \texttt{\tiny{[GM$|$GM]}}
                \texttt{game\_result = LOSE} \\
            }
        }
    }
    & & \\ \\

\end{supertabular}
}

\end{document}
